\begin{tablalafrox}{Diccionario del Producto 7: Innovaciones verificadas. Fuente: según cita por cuenta en cada encabezado.}{tab:T14-dic-prod-7}
  {L{1.2cm} Y Y Y C{0.9cm} L{2.3cm}}
  {\cab{Código} & \cab{Paquete de trabajo} & \cab{Entregable} & \cab{Criterio de aceptación} & \cab{Resp.} & \cab{Período}}
  \multicolumn{6}{>{\columncolor{lafroxBg}\RaggedRight\arraybackslash}p{\dimexpr\textwidth-2\tabcolsep\relax}}{\textbf{Cuenta 7.1 — Innovación 1: seguimiento de vencimiento posentrega en el local.} Esta cuenta reúne el catálogo de acciones, el estimador de saldo y ritmo de reposición, la lista sin conexión y el piloto para el seguimiento del vencimiento posentrega. Sus cuatro paquetes se ejecutan en la Etapa 1, incluida la marcha blanca. {\color{lafroxGris}Fuente: (Bases Administrativas, arts. 17.2–17.6 y Formulario T-14, p. 63); (Bases Técnicas Transversales, secciones 17.2–17.6, pp. 32–34).}} \tabularnewline*
  7.1.1 & Innovación 1: documento de catálogo de acciones y protocolo del piloto (INN-01.P1) & Documento de catálogo de acciones y protocolo del piloto. & Aprobado por Comercial y Calidad. & DAT & E1 · Meses 7 a 9 \\
  7.1.2 & Innovación 1: estimador de saldo y ritmo de reposición (INN-01.P2) & Estimador de saldo y ritmo de reposición. & Supera QA y cierra antes del límite contractual del H4 (30-11-2027), sin integrar su entrega ni su acta. & DAT & E1 · Meses 9 y 10 (límite del H4) \\
  7.1.3 & Innovación 1: función de lista sin conexión en la aplicación de preventa (INN-01.P3) & Función de lista sin conexión en la aplicación de preventa. & Funciona sin señal y cierra antes del límite contractual del H4 (30-11-2027), sin integrar su entrega ni su acta. & DAT & E1 · Meses 9 y 10 (límite del H4) \\
  7.1.4 & Innovación 1: informe del piloto en marcha blanca y habilitación (INN-01.P4) & Informe del piloto en marcha blanca y habilitación. & Al mes 15, al menos el 70 \% de los avisos confirma el riesgo, con al menos 100 avisos. & DAT & E1 · Meses 11 a 15 (H5 y marcha blanca) \\
  \midrule
  \multicolumn{6}{>{\columncolor{lafroxBg}\RaggedRight\arraybackslash}p{\dimexpr\textwidth-2\tabcolsep\relax}}{\textbf{Cuenta 7.2 — Innovación 2: reproducción de incidentes de terreno.} Esta cuenta reúne la captura y reproducción de evidencia de incidentes de terreno y su validación comparada durante la implementación. En Operación se mantienen los escenarios de regresión mediante informes trimestrales, desde el mes 21. {\color{lafroxGris}Fuente: (Bases Administrativas, arts. 17.2–17.6 y Formulario T-14, p. 63); (Bases Técnicas Transversales, secciones 17.2–17.6, pp. 32–34).}} \tabularnewline*
  7.2.1 & Innovación 2: informe de medición inicial y contrato de evidencia (INN-02.P1) & Informe de medición inicial y contrato de evidencia. & Aprobado por el Encargado de Seguridad. & CAL & E1 · Meses 2 y 3 \\
  7.2.2 & Innovación 2: circuito mínimo de reproducción en la preventa (INN-02.P2) & Circuito mínimo de reproducción en la preventa. & Un incidente se reproduce en QA a partir de su evidencia. & CAL & E1 · Meses 3 a 5 \\
  7.2.3 & Innovación 2: captura de evidencia en las aplicaciones de preventa, preparación y reparto (INN-02.P3) & Captura de evidencia en las aplicaciones de preventa, preparación y reparto. & Las tres aplicaciones envían evidencia válida antes del H4. & CAL & E1 · Meses 5 a 10 (antes del H4) \\
  7.2.4 & Innovación 2: informe de validación comparada (INN-02.P4, parte 1) & Informe de validación comparada durante la implementación. & Al mes 15 se cumplen dos metas: el tiempo de diagnóstico baja 30 \%, con al menos 20 casos comparables, y al menos el 80 \% de los incidentes se reproduce. & CAL & E1 · Meses 11 a 15 (validación en meses 13 a 15) \\
  7.2.5 & Innovación 2: mantenimiento de escenarios de regresión (INN-02.P4, segunda parte) & Informe trimestral. & Ningún incidente se reabre por la misma causa, y ninguna evidencia contiene datos personales. & CAL & Operación · Desde el mes 21. \\
  \midrule
  \multicolumn{6}{>{\columncolor{lafroxBg}\RaggedRight\arraybackslash}p{\dimexpr\textwidth-2\tabcolsep\relax}}{\textbf{Cuenta 7.3 — Innovación 3: vida útil remanente por historia térmica del lote.} Esta cuenta reúne los parámetros, la asociación lote--sensor, la vista de vida útil remanente y su validación con Calidad en marcha blanca. Durante la Operación se revisan periódicamente los parámetros con aprobación del Líder de Calidad de LafroX y de Calidad del CLIENTE. {\color{lafroxGris}Fuente: (Bases Administrativas, arts. 17.2–17.6 y Formulario T-14, p. 63); (Bases Técnicas Transversales, secciones 17.2–17.6, pp. 32–34).}} \tabularnewline*
  7.3.1 & Innovación 3: documento de familias priorizadas y parámetros aprobados (INN-03.P1) & Documento de familias priorizadas y parámetros aprobados. & Los parámetros están aprobados por el Líder de Calidad y Pruebas y Calidad del CLIENTE. & DAT & E1 · Meses 3 a 5 \\
  7.3.2 & Innovación 3: asociación lote–sensor en el borde y en el camión (INN-03.P2) & Asociación lote–sensor en el borde y en el camión. & Un lote recorre todo su camino con historia completa en QA. & DAT & E1 · Meses 5 a 9 \\
  7.3.3 & Innovación 3: vista en M9 y orden de salida sugerido (INN-03.P3) & Vista en M9 y orden de salida sugerido. & No cambia la regla de Calidad y cierra antes del límite contractual del H4 (30-11-2027), sin integrar su entrega ni su acta. & DAT & E1 · Meses 9 y 10 (límite del H4) \\
  7.3.4 & Innovación 3: informe de validación en marcha blanca con Calidad (INN-03.P4) & Informe de validación en marcha blanca con Calidad. & Al mes 15 se cumplen dos metas: al menos el 95 \% de los lotes tiene historia completa, y al menos el 80 \% de las estimaciones queda dentro de ±15 \%. & DAT & E1 · Meses 11 a 15 (H5 y marcha blanca) \\
  7.3.5 & Innovación 3: informes de revisión periódica de parámetros (INN-03.P5) & Parámetros ajustados con aprobación del Líder de Calidad de LafroX y de Calidad del CLIENTE. & Cada revisión periódica queda documentada y aprobada por el Líder de Calidad de LafroX y Calidad del CLIENTE. Si la exactitud baja del 80 \% dos veces seguidas, se hace un ajuste extraordinario. & SRE & Operación · Meses 24, 27, 30, 36, 39, 42, 48, 51 y 54 del proyecto. \\
  \midrule
  \multicolumn{6}{>{\columncolor{lafroxBg}\RaggedRight\arraybackslash}p{\dimexpr\textwidth-2\tabcolsep\relax}}{\textbf{Cuenta 7.4 — Innovación 4: tramo variable vinculado a resultados.} Esta cuenta reúne la fórmula contractual del tramo variable vinculado a resultados durante la Etapa 2. La Operación comprende la línea base, tres liquidaciones en sombra y las 33 liquidaciones mensuales de los meses 24 a 56. {\color{lafroxGris}Fuente: (Bases Administrativas, arts. 17.2–17.6 y Formulario T-14, p. 63); (Bases Técnicas Transversales, secciones 17.2–17.6, pp. 32–34).}} \tabularnewline*
  7.4.1 & Innovación 4: fórmula aprobada del tramo variable (INN-04.P1) & Fórmula aprobada del tramo variable (reglas de atribución, canasta e índices). & Finanzas del CLIENTE aprueba al cierre del paquete la fórmula inicial, simulada con casos. & JP & E2 · Ejecución nominal: 1 al 9 de junio de 2028 (mes 17). Ventana de asignación: meses 17 a 20. \\
  7.4.2 & Innovación 4: línea base y liquidaciones en sombra (INN-04.P2) & Línea base y tres liquidaciones de prueba. & Finanzas reproduce las tres sin diferencias. & JP & Operación · Meses 21 a 23. \\
  7.4.3 & Innovación 4: liquidaciones mensuales del tramo variable (INN-04.P3) & Las 33 liquidaciones mensuales (meses 24 a 56) y el registro de las recomendaciones. & Finanzas reproduce cada liquidación, y cada recomendación tiene su respuesta registrada. & JP & Operación · Mensual, desde el mes 24. \\
  \midrule
  \multicolumn{6}{>{\columncolor{lafroxBg}\RaggedRight\arraybackslash}p{\dimexpr\textwidth-2\tabcolsep\relax}}{\textbf{Cuenta 7.5 — Innovación 5: hoja de negocio del almacenero.} Esta cuenta reúne el prototipo, la evaluación de privacidad, el cálculo y la certificación de la hoja de negocio del almacenero durante la Etapa 2. En Operación considera su despliegue gradual por rutas y el bloque 4 condicionado al dictamen legal. {\color{lafroxGris}Fuente: (Bases Administrativas, arts. 17.2–17.6 y Formulario T-14, p. 63); (Bases Técnicas Transversales, secciones 17.2–17.6, pp. 32–34).}} \tabularnewline*
  7.5.1 & Innovación 5: prototipo e informe de la prueba de comprensión (INN-05.P1) & Prototipo e informe de la prueba de comprensión. & Al menos el 80 \% de los 24 almaceneros entiende la hoja. & IMP & E2 · Meses 13 y 14 \\
  7.5.2 & Innovación 5: evaluación de privacidad y diseño de la impresión en cabina (INN-05.P2) & Evaluación de privacidad y diseño de la impresión en cabina. & Aprobada por el Encargado de Seguridad, según la Ley 21.719. & IMP & E2 · Meses 14 y 15 \\
  7.5.3 & Innovación 5: cálculo y bloques 1 a 3 de la hoja de negocio (INN-05.P3) & Cálculo y bloques 1 a 3 de la hoja de negocio. & Entra al H9 y funciona sin señal. & IMP & E2 · Meses 15 a 17 (H9) \\
  7.5.4 & Innovación 5: hoja certificada, en marcha blanca y en producción (INN-05.P4, primera parte) & Hoja certificada, en marcha blanca y en producción. & Incluida en el acta del H12. & IMP & E2 · Meses 18 a 21 (H10 a H12) \\
  7.5.5 & Innovación 5: despliegue gradual por rutas (INN-05.P4, segunda parte) & Hoja disponible en todas las rutas. & Al mes 27, la hoja se muestra en al menos el 90 \% de los locales que la aceptan, y al menos el 85 \% de las copias pedidas se entrega dentro de dos visitas. & SRE & Operación · Meses 22 a 26. \\
  7.5.6 & Innovación 5: bloque 4 condicionado (INN-05.P5) & Dictamen legal y, si es favorable, el bloque activado. & No se activa sin dictamen legal favorable y aprobación de privacidad por el Encargado de Seguridad; todas las emisiones cumplen el umbral de agrupación. & SRE & Operación · Meses 24 a 27. \\
\end{tablalafrox}
